\NativeSlot{48}{44}{%
% p066_s01

\ExamQuestion{5.}{5}%
奥运五环连在一起代表着五大洲的人们友好相处，齐聚奥运盛会。它可以由\mbox{（\hspace{1.8em}）}形（平移、旋转）（在正确答案下面画“√”）而成，它有\mbox{（\hspace{1.8em}）}条对称轴，请在图中画出来。\par
\BookFigure{figures/p066_s01-01.tex}
}
\NativeSlot{128}{45}{%
\par\noindent%
\ExamTile{.485}{%
% p066_s02

\ExamQuestion{6.}{5}%
找规律，画一画。\par\noindent （1）在空白方格里补画图形。\par
\BookFigure{figures/p066_s02-01.tex}
}\hfill%
\ExamTile{.485}{%
% p066_s03

\ExamQuestion{（2）}{5}%
在第四格补画图形。\par
\BookFigure{figures/p066_s03-01.tex}
}\par
}
\NativeSlot{211}{125}{%
\par\noindent%
\ExamTile{.54}{%
% p066_s04

\ExamQuestion{7（①）}{5}%
要将图二恢复到图一，将①部分先向上\mbox{（\hspace{1.8em}）}\mbox{（\hspace{1.8em}）}格，再向\mbox{（\hspace{1.8em}）}\mbox{（\hspace{1.8em}）}\mbox{（\hspace{1.8em}）}格。图一、图二为4×4方格拼马图。\par

% p066_s05

\ExamQuestion{7（②）}{5}%
要将图二恢复到图一，将②部分绕\mbox{（\hspace{1.8em}）}点\mbox{（\hspace{1.8em}）}时针旋转\mbox{（\hspace{1.8em}）}度，再向\mbox{（\hspace{1.8em}）}平移\mbox{（\hspace{1.8em}）}格。图二右下方方格的左下顶点为A，右上顶点为B。\par
}\hfill%
\ExamTile{.435}{%
\BookFigure{figures/p066_s04-01.tex}
}\par
}
\NativeSlot{418.25}{74}{%
\ExamGroup{二、选择（将正确答案的序号填在括号里）。{\normalfont\kaishu\fontsize{10.8bp}{14bp}\selectfont （每空2分，共14分）}}

% p066_s06

\ExamQuestion{1.}{5}%
下面运动现象中，属于平移运动的有\mbox{（\hspace{1.8em}）}个。图示依次是抽拉键盘托架、订书机、旋转闸、门环、抽屉。\par
\ExamOptions{A. \(1\)}{B. \(2\)}{C. \(3\)}{D. \(4\)}
\BookFigure{figures/p066_s06-01.tex}
}
\NativeSlot{543}{89}{%
% p066_s07

\ExamQuestion{2.}{5}%
下面几组图形中，甲图旋转得到乙图的是\mbox{（\hspace{1.8em}）}。\par
\ExamOptions{A. 一组}{B. 二组}{C. 一组和二组}{D. 二组和三组}
\BookFigure{figures/p066_s07-01.tex}
}
\NativeSlot{693}{58}{%
% p066_s08

\ExamQuestion{3.}{5}%
左图黑白风车图案是用\mbox{（\hspace{1.8em}）}图形，\mbox{（\hspace{1.8em}）}得到的。\par
\par\noindent%
\ExamTile{.24}{%
\BookFigure{figures/p066_s08-01.tex}
}\hfill%
\ExamTile{.73}{%
\BookFigure{figures/p066_s08-02.tex}
}\par
}
